\section*{РОЗДІЛ 4. ДО ТЕХНІЧНИХ МЕХАНІЗМІВ ІНТЕГРАЦІЇ GENAI В ІНФОРМАЦІЙНІ СИСТЕМИ ЗАКЛАДІВ ВИЩОЇ ОСВІТИ}

\subsection*{4.1. Системи CRIS та стандарти інтероперабельності}

\subsubsection*{4.1.1. Огляд провідних CRIS-платформ}

\textbf{Current Research Information Systems (CRIS)} доцільно трактувати як інструментальний клас систем управління дослідницькою інформацією, що забезпечує консолідацію, актуалізацію та повторне використання даних про наукову діяльність організації в єдиному інформаційному контурі. На відміну від інституційних репозитаріїв, орієнтованих передусім на зберігання повних текстів, CRIS підтримують ширший спектр сутностей і зв'язків (публікації, проєкти, гранти, обладнання, співпраці), що робить їх придатними для формування управлінської та аналітичної звітності на рівні підрозділів і всієї установи \cite{Bryant2017RIMDefining}.

Екосистема CRIS охоплює як комерційні, так і відкриті реалізації. Комерційні платформи \textbf{Pure (Elsevier)} та \textbf{Converis (Clarivate)} репрезентують enterprise-підхід: вони поєднують модулі профілів дослідників, робочі процеси (workflow) та інтеграції з наукометричними індексами й реєстрами ідентифікаторів (зокрема ORCID) для автоматизованого наповнення й верифікації метаданих \cite{ORCID2024APIDocs}. Натомість \textbf{VIVO} позиціонується як відкрита семантична платформа, що спирається на онтологічне подання даних (RDF/OWL) та підтримує підхід Linked Data; це підвищує інтероперабельність і спрощує адаптацію моделі даних під локальні потреби.

Окрему групу становлять рішення, що еволюціонували з репозитарних платформ у бік CRIS-функціональності. У цьому контексті \textbf{DSpace-CRIS} можна розглядати як приклад поєднання репозитарію та CRIS на основі вимог до інтероперабельності, зокрема узгодження метаданих і процедур обміну з екосистемою OpenAIRE \cite{OpenAIRE2023GuidelinesCRIS}.

\subsubsection*{4.1.2. Стандарт CERIF та OpenAIRE Guidelines}

\textbf{CERIF (Common European Research Information Format)} є одним із базових європейських стандартів для формалізованого подання та обміну дослідницькою інформацією між організаціями й інформаційними системами; його розвиток і підтримка здійснюються консорціумом euroCRIS \cite{euroCRIS2024CERIF}. У логіці інтеграції даних CERIF задає узгоджену сутнісно-зв'язкову модель, у межах якої ключові об'єкти дослідницької діяльності (дослідник, організаційна одиниця, проєкт, публікація, результат/продукт, фінансування тощо) описуються стандартизовано та пов'язуються типізованими відношеннями. Семантичну узгодженість забезпечують класифікаційні схеми (classification schemes), що уніфікують типи результатів, ролі та статуси й тим самим зменшують неоднозначність при міжсистемному обміні.

У практичних сценаріях інтероперабельності вагому роль відіграють \textbf{OpenAIRE Guidelines for CRIS Managers v1.2.0}, які задають профіль застосування CERIF для забезпечення сумісності з European Open Science Cloud (EOSC) \cite{OpenAIRE2023GuidelinesCRIS}. З погляду технічної інтеграції ці настанови конкретизують протокольний рівень обміну (насамперед OAI-PMH), мінімально достатній набір обов'язкових та рекомендованих полів для кожної сутності, а також правила валідації та контролю якості метаданих.

\subsection*{4.2. Model Context Protocol (MCP) для інтеграції ШІ-інструментів}

\subsubsection*{4.2.1. Архітектура та принципи MCP}

\textbf{Model Context Protocol (MCP)} запропоновано як відкритий протокол, покликаний стандартизувати взаємодію LLM-застосунків із зовнішніми джерелами даних та інструментами, тобто забезпечити переносимий механізм підключення контексту та дій до моделі \cite{Anthropic2024MCP}. Його ключовою мотивацією є усунення проблеми $N\times M$ інтеграцій: у традиційній схемі кожну комбінацію «LLM~--- сервіс» необхідно реалізовувати окремим конектором, тоді як MCP фіксує уніфікований контракт, який реалізується один раз на стороні клієнта та один раз на стороні сервера \cite{Anthropic2024MCP}.

На технічному рівні MCP визначає формат повідомлень і транспорт (зокрема JSON-RPC 2.0 через stdio для локальних процесів або мережеві режими на кшталт HTTP/SSE), а також передбачає типові механізми захисту доступу (наприклад OAuth 2.0) і семантичне версіонування для підтримки сумісності між реалізаціями \cite{Anthropic2024MCP}. Концептуально протокол оперує трьома базовими абстракціями: \textit{tools} (формалізовані дії), \textit{resources} (читані джерела даних) та \textit{prompts} (повторно використовувані шаблони взаємодії).

\subsubsection*{4.2.2. Клієнт-серверна архітектура}

Архітектура MCP передбачає розподіл ролей між \textbf{MCP-клієнтом}, який інкапсулює LLM та логіку оркестрації, і \textbf{MCP-сервером}, що надає доступ до конкретного доменного ресурсу (файлова система, СУБД, зовнішні API або корпоративні системи на кшталт CRIS). Типовий життєвий цикл включає виявлення доступних можливостей (discovery), вибір релевантних інструментів, виклик (invocation) та повернення результатів із подальшою інтеграцією їх у відповідь або робочий процес користувача \cite{Anthropic2024MCP}.

\subsubsection*{4.2.3. Доступні SDK та готові сервери}

Практична цінність MCP полягає у можливості швидкого створення серверів до різних джерел даних завдяки наявності SDK та еталонних реалізацій. У прикладних сценаріях наукового адміністрування найбільш релевантними є конектори до ідентифікаційних і бібліометричних сервісів (ORCID, OpenAlex), а також до локальних сховищ і СУБД; це дозволяє GenAI-агенту відтворювано отримувати метадані, здійснювати верифікацію та агрегувати показники за узгодженими правилами [14, 15, 34].

\subsection*{4.3. Підхід «Golden Record» та алгоритм злиття даних}

\subsubsection*{4.3.1. Концепція Golden Record}

Інтеграція даних із різнорідних наукометричних та адміністративних джерел неминуче породжує дублікати й конфлікти атрибутів для одних і тих самих об'єктів (публікацій, авторів, проєктів). У таких умовах концепція \textbf{Golden Record} (Master Data, Single Source of Truth) задає методологічну рамку, в межах якої формується єдиний «авторитетний» запис, що агрегує найбільш повні та узгоджені значення з доступних джерел і водночас забезпечує трасованість (provenance) прийнятих рішень.

Для публікацій Golden Record зазвичай спирається на стабільний унікальний ідентифікатор (передусім DOI), нормалізовані метадані (назва, авторський склад, джерело, рік) та агреговані показники впливу (наприклад, узгоджене або максимальне значення цитувань за різними індексами). Принципово важливою є фіксація провенансу: які джерела надали певні атрибути та за якими правилами ці атрибути було обрано або об'єднано.

\subsubsection*{4.3.2. Стратегії визначення «переможця»}

Розв'язання конфліктів при злитті даних можна трактувати як задачу вибору «переможця» для кожного атрибута. На практиці доцільно комбінувати кілька комплементарних стратегій. \textit{Ручний відбір} забезпечує максимальну точність, однак не масштабується на великі масиви даних. \textit{Пріоритезація джерел} задає ієрархію довіри та дозволяє автоматизувати вибір, якщо відомо, які системи є первинними для конкретних полів. У наукометричному контексті обґрунтовано надавати вищий пріоритет реєстрам ідентифікаторів та агрегаторам метаданих (наприклад Crossref, OpenAlex), тоді як дані самозвіту (ORCID) варто розглядати як важливе, але потенційно неповне доповнення [14, 15]. \textit{Вибір найновішого} є корисним для атрибутів, чутливих до часу оновлення. Нарешті, \textit{атрибутна агрегація} дозволяє застосовувати найадекватнішу стратегію окремо для кожного поля (наприклад, для цитувань --- агрегування/максимум; для ключових слів --- об'єднання без повторів).

\subsubsection*{4.3.3. Розроблений алгоритм Golden Record}

У межах запропонованого підходу початковим кроком є \textbf{нормалізація DOI}: приведення до нижнього регістру, вилучення URL-префіксів (на кшталт \texttt{https://doi.org/}), усунення зайвих пробілів та перевірка відповідності формату. Така уніфікація зменшує кількість хибних розбіжностей і підвищує точність первинного зіставлення записів у різних джерелах.

\textit{Псевдокод для кроків 1--5 (\texttt{normalize\_doi}, \texttt{primary\_match}, \texttt{secondary\_match}, \texttt{normalize\_title}, \texttt{merge} та структури провенансу) наведено у Додатку A (лістинги~\ref{lst:a-normalize-doi}--\ref{lst:a-provenance}).}

\subsubsection*{4.3.4. Референсна система OpenAIRE GDup}

Як приклад промислово придатного підходу до дедуплікації в інфраструктурах відкритої науки доцільно розглянути \textbf{OpenAIRE GDup}, орієнтовану на обробку великих графів наукових даних \cite{OpenAIRE2023GuidelinesCRIS}. Архітектурно система поєднує фазу \textit{блокування (blocking)} (групування записів за ключами блокування для уникнення повного попарного порівняння), подальші порівняння в межах блоків із використанням метрик схожості (наприклад Jaccard, Levenshtein, TF--IDF cosine) та контур керування \textit{ground truth}, у межах якого ручні позначки кураторів використовуються для корекції та поступового вдосконалення правил зіставлення \cite{OpenAIRE2023GuidelinesCRIS}.

\subsection*{4.4. Референсна архітектура інтеграції GenAI + CRIS}

\begin{figure}[ht]
\centering
\begin{tikzpicture}[
  font=\small,
  node distance=6mm,
  box/.style={draw, rounded corners, align=center, inner sep=6pt, minimum width=0.88\linewidth}
]
  \node[box, fill=gray!10] (llm) {LLM Service Layer\\\footnotesize оркестрація інструментів \;|\; керування контекстом \;|\; генерація звітів/відповідей};

  \node[box, below=of llm, fill=blue!6] (mcp) {MCP Server Layer\\\footnotesize \texttt{mcp-orcid} \;|\; \texttt{mcp-openalex} \;|\; \texttt{mcp-cris} \;|\; \texttt{mcp-local}};

  \node[box, below=of mcp, fill=green!6] (fusion) {Data Fusion Layer\\\footnotesize disambiguation \;|\; Golden Record Assembly \;|\; linking (DOI/ORCID/ROR)};

  \node[box, below=of fusion, fill=orange!10] (api) {API Gateway Layer\\\footnotesize OAuth 2.0 \;|\; rate limiting \;|\; протокольна трансляція (REST/GraphQL/SOAP)};

  \node[box, below=of api, fill=gray!5] (sources) {Джерела даних та системи\\\footnotesize ORCID \;|\; OpenAlex \;|\; Crossref/Scopus/WoS \;|\; локальні БД/файли \;|\; CRIS};

  \draw[-{Stealth[length=3mm]}, thick] (sources) -- (api);
  \draw[-{Stealth[length=3mm]}, thick] (api) -- (fusion);
  \draw[-{Stealth[length=3mm]}, thick] (fusion) -- (mcp);
  \draw[-{Stealth[length=3mm]}, thick] (mcp) -- (llm);
\end{tikzpicture}
\caption{Структурна схема референсної архітектури інтеграції GenAI + CRIS}
\label{fig:genai-cris-arch}
\end{figure}